\documentclass[a4paper,11pt]{ltjsarticle}
\usepackage[haranoaji]{luatexja-preset}
\usepackage{amsmath,amssymb}
\usepackage[top=12mm,bottom=10mm,left=14mm,right=14mm]{geometry}
\usepackage{array}
\usepackage{tikz}
\usetikzlibrary{calc}
\pagestyle{empty}
\setlength{\parindent}{0pt}
\newlength{\qcol}\setlength{\qcol}{0.47\textwidth}
\newlength{\qgap}
\newlength{\anscolw}\setlength{\anscolw}{0.30\textwidth}
\newlength{\ansh}
\newcommand{\Q}[2]{\parbox[t]{\qcol}{\raggedright (#1)\hspace{0.6em}$\displaystyle #2$}}
\newcommand{\QR}[4]{\Q{#1}{#2}\hfill\Q{#3}{#4}\par\vspace{\qgap}}
\newcommand{\anscell}[1]{\rule[-\ansdrop]{0pt}{\ansh}(#1)}
\newlength{\ansdrop}

\begin{document}
{\large\bfseries 計算問題　6}\hfill 名前(\underline{\hspace{32mm}})\par\vspace{3mm}
■（1）〜（16）の計算をしなさい。（17）〜（20）は方程式を解きなさい。\par\vspace{3mm}
\setlength{\qgap}{5.4mm}
\QR{1}{-9+(-3)}{2}{(-5)-(+12)}
\QR{3}{(3a+8)+(2a-5)}{4}{(4x+2)-(2x-6)}
\QR{5}{2x\times(-5)}{6}{-40a\div8}
\QR{7}{(+10)+(-5)+(-15)}{8}{(-54)\div(+6)}
\QR{9}{(-12)\times(-13)}{10}{\dfrac{1}{2}(6a+4)-\dfrac{1}{3}(9a-12)}
\QR{11}{\dfrac{x+3}{4}\times8}{12}{(18a-6)\div3}
\QR{13}{9\div12\times(-16)}{14}{5+2^3+(-4)}
\QR{15}{\dfrac{1}{2}(2x+4)+\dfrac{1}{3}(3x+6)}{16}{3-(-9)\div\left(-\dfrac{3}{4}\right)}
\QR{17}{5x+1=-9}{18}{3(x+4)=2x+9}
\QR{19}{0.3x+0.7=4}{20}{\dfrac{1}{3}x=\dfrac{1}{5}x+2}
\vspace{1mm}
\Q{21}{1+2+3+4+5+6+7+8+9+10}\par\vspace{3mm}
\setlength{\ansh}{9.5mm}\setlength{\ansdrop}{6mm}%
\begin{center}\begin{tabular}{|p{\anscolw}|p{\anscolw}|p{\anscolw}|}\hline
\anscell{1} & \anscell{2} & \anscell{3} \\\hline
\anscell{4} & \anscell{5} & \anscell{6} \\\hline
\anscell{7} & \anscell{8} & \anscell{9} \\\hline
\anscell{10} & \anscell{11} & \anscell{12} \\\hline
\anscell{13} & \anscell{14} & \anscell{15} \\\hline
\anscell{16} & \anscell{17} & \anscell{18} \\\hline
\anscell{19} & \anscell{20} & \anscell{21} \\\hline
\end{tabular}\end{center}

\end{document}
